% Generated from locale/tr/content/many-valued-logic/syntax-and-semantics/formulas.tex; do not hand-edit.


\olfileid[tr]{mvl}{syn}{fml}

\olsection{\usetoken{P}{formula}}

\begin{defn}[Formül]
\ollabel{defn:formulas}
Bir $\Lang L$ önerme dilinin $\Frm[L]$ \emph{!!{formula}s} kümesi
tümevarımlı olarak şöyle tanımlanır:
\begin{enumerate}
\item Her !!{propositional variable} $\Obj p_i$, atomik bir !!a{formula}
  oluşturur.
\item $\Lang L$ dilinin her $0$-konumlu bağlacı (önerme sabiti), atomik
  bir !!a{formula} oluşturur.
\item $\star$, $\Lang L$ dilinin $n$-konumlu bir bağlacı ve $!A_1$,
\dots, $!A_n$ !!{formula}s ise $\star(!A_1,\dots,!A_n)$ bir
!!a{formula} olur.
\tagitem{limitClause}{Başka hiçbir şey !!a{formula} değildir.}{}
\end{enumerate}
$\star$ tek konumluysa $\star(!A_1)$ çoğu zaman yalnızca $\star !A_1$;
$\star$ iki konumluysa $\star(!A_1,!A_2)$ çoğu zaman
$(!A_1\star !A_2)$ biçiminde yazılır.
\end{defn}

Her zamanki gibi çoğu zaman en dıştaki parantezleri sessizce atacağız.

\begin{ex}
Standart $\Lang{L_0}$ dilinde
$\Obj p_1 \lif (\Obj p_1 \land \lnot \Obj p_2)$ bir formüldür. Çarpım
mantığının dilinde bunun yerine
$\Obj p_1 \lif (\Obj p_1 \odot \lnot \Obj p_2)$ yazılırdı. Dile tek
konumlu $\triangle$ eklersek
$\triangle (\Obj p_1 \land \Obj p_2) \lif
(\triangle \Obj p_1 \land \triangle \Obj p_2)$ gibi formüllerimiz de olur.
\end{ex}

